\documentclass[11pt]{article}
\usepackage[margin=1in]{geometry}
\usepackage{amsmath,amssymb,amsthm}
\usepackage{graphicx}
\usepackage[hidelinks]{hyperref}

\newtheorem{theorem}{Theorem}
\newtheorem{lemma}{Lemma}
\newtheorem{remark}{Remark}

\newcommand{\N}{\mathbb N}
\newcommand{\K}{\mathbb K}
\newcommand{\eps}{\varepsilon}

\title{A Non-Absolute Schauder Frame in a Nuclear Frechet Space}
\author{fa\_banach\_001, agent\_lane\_16}
\date{2026-07-03}

\begin{document}
\maketitle

\section*{Source Question}

This packet concerns Remark 2.6 of Cassese--Ramos, \emph{On the regularity of Fourier interpolation formulas}, arXiv:2503.15733. The paper defines a Schauder frame in a locally convex space \(X\) as a sequence \(((x_n,y_n))_n\), with \(x_n\in X\) and \(y_n\in X'\), such that
\[
        x=\sum_{n=1}^{\infty} y_n(x)x_n
\]
for every \(x\in X\), with convergence in the topology of \(X\). Immediately after this definition, the paper asks whether, in view of the corresponding theorem for bases, a Schauder frame in a nuclear Frechet space must be absolute.

\begin{center}
\includegraphics[width=0.98\textwidth]{figures/open_problem_crop.png}
\end{center}

We interpret ``absolute Schauder frame'' in the direct frame analogue of the source paper's Definition 2.1 for absolute bases: for every \(x\in X\) and every continuous seminorm \(p\), the scalar series
\[
        \sum_{n=1}^{\infty} |y_n(x)|\,p(x_n)
\]
converges. The counterexample below also fails the stronger bounded-set version obtained by replacing \(|y_n(x)|\) with \(p_B(y_n)=\sup_{x\in \operatorname{co}(B)} |y_n(x)|\), by taking \(B=\{e_1\}\).

\section*{Answer}

The answer is negative. There is a Schauder frame in the nuclear Frechet space \(s\) of rapidly decreasing scalar sequences whose reconstruction series converges in \(s\) for every vector but is not absolute even at one basis vector.

\section*{Proof Intuition}

Start with the usual absolute basis \((e_j)\) of \(s\). Replace each single basis coefficient by a full conditionally convergent scalar row: choose scalars \(r_k\) with \(\sum_k r_k=1\) but \(\sum_k |r_k|=\infty\). The \(j\)-th coordinate is then reconstructed as \(a_j\sum_k r_k e_j=a_j e_j\). Enumerating the double sequence by diagonals makes every fixed coordinate see longer and longer initial sums of \(\sum r_k\), while the partial sums remain uniformly bounded. Rapid decay of \(a_j\) controls the moving tail over \(j\). Thus the frame converges in \(s\), but absolute convergence fails on \(e_1\) because it contains the divergent series \(\sum_k |r_k|\).

\section*{The Counterexample}

Let
\[
    s=\left\{a=(a_j)_{j\ge 1}:\ p_m(a):=\sup_{j\ge 1} j^m |a_j|<\infty
    \text{ for every } m\in\N_0\right\}.
\]
This is the usual Frechet space of rapidly decreasing sequences. It is nuclear: indeed, for \(m\in\N_0\) choose \(\ell>m+1\). The canonical map from the Banach completion for \(p_\ell\) to the Banach completion for \(p_m\) has the representation
\[
        a\mapsto \sum_{j=1}^{\infty} e_j^*(a)e_j,
\]
where \(\|e_j^*\|_{p_\ell}\le j^{-\ell}\) and \(p_m(e_j)=j^m\), so
\[
        \sum_{j=1}^{\infty} \|e_j^*\|_{p_\ell}\,p_m(e_j)
        \le \sum_{j=1}^{\infty} j^{m-\ell}<\infty.
\]
This gives the standard nuclearity criterion for \(s\).

Set
\[
        r_k=\frac{(-1)^{k+1}}{k\log 2},\qquad k\ge 1.
\]
Then \(\sum_{k=1}^{\infty} r_k=1\), \(\sum_k |r_k|=\infty\), and the partial sums
\(\sigma_K=\sum_{k=1}^K r_k\), with \(\sigma_0=0\), are bounded. Let
\[
        C_0=\sup_{K\ge 0} |\sigma_K|<\infty.
\]

Enumerate \(\N\times\N\) by diagonals, in increasing order of \(j+k\), and within each diagonal in increasing order of \(j\). Write the resulting enumeration as
\[
        n\mapsto (j_n,k_n).
\]
Define
\[
        x_n=e_{j_n}\in s,\qquad
        y_n=r_{k_n} e_{j_n}^*\in s',
\]
where \(e_j^*(a)=a_j\). Each \(y_n\) is continuous.

\begin{theorem}
The sequence \(((x_n,y_n))_{n\ge 1}\) is a Schauder frame for the nuclear Frechet space \(s\), but it is not an absolute Schauder frame.
\end{theorem}

\begin{proof}
Fix \(a=(a_j)\in s\), and let
\[
        S_N a=\sum_{n=1}^N y_n(a)x_n.
\]
Because of the diagonal ordering, for each \(j\) the terms from the \(j\)-th row appear in their natural order \(k=1,2,\ldots\). Hence
\[
        S_Na=\sum_{j=1}^{\infty} a_j \sigma_{K_j(N)} e_j,
\]
where \(K_j(N)\in\N_0\), \(K_j(N)\to\infty\) for each fixed \(j\), and
\[
        \sup_{j,N} |\sigma_{K_j(N)}|\le C_0.
\]
The displayed sum is finite for each \(N\), but this coordinate formula is the convenient way to compare it with \(a\).

We prove \(S_Na\to a\) in every seminorm \(p_m\). Let \(\eta>0\). Since \(a\in s\), choose \(J\) so large that
\[
        \sup_{j>J} j^m |a_j| < \frac{\eta}{2(C_0+1)}.
\]
For the finitely many indices \(1\le j\le J\), the convergence
\(\sigma_{K_j(N)}\to 1\) implies that, for all sufficiently large \(N\),
\[
        \max_{1\le j\le J} j^m |a_j|\,|\sigma_{K_j(N)}-1|<\frac{\eta}{2}.
\]
For these \(N\),
\[
\begin{aligned}
        p_m(S_Na-a)
        &=\sup_{j\ge 1} j^m |a_j|\,|\sigma_{K_j(N)}-1|\\
        &\le \max_{1\le j\le J} j^m |a_j|\,|\sigma_{K_j(N)}-1|
          +(C_0+1)\sup_{j>J} j^m |a_j|\\
        &<\eta.
\end{aligned}
\]
Thus \(S_Na\to a\) in \(s\), so \(((x_n,y_n))\) is a Schauder frame.

It is not absolute. Take \(a=e_1\) and \(p_0(a)=\sup_j |a_j|\). The terms with \(j_n=1\) contribute
\[
        |y_n(e_1)|\,p_0(x_n)=|r_{k_n}|\,p_0(e_1)=|r_{k_n}|,
\]
and therefore
\[
        \sum_{n=1}^{\infty} |y_n(e_1)|\,p_0(x_n)
        \ge \sum_{k=1}^{\infty} |r_k|
        =\infty.
\]
Consequently the frame is not absolute. The same computation with \(B=\{e_1\}\) shows failure of the stronger bounded-set absolute condition as well.
\end{proof}

\section*{Verification Notes}

The proof is purely analytical. The only external source used for the target statement and definitions is the Cassese--Ramos paper. A finite numerical sanity check in \texttt{code/check\_diagonal\_frame.py} illustrates the diagonal partial-sum mechanism and the growth of \(\sum_{k\le N}|r_k|\); it is not used as proof.

\section*{Novelty Check}

Before promotion, the run indexes were searched for close variants of
\[
\texttt{Schauder frame},\quad \texttt{nuclear Frechet},\quad
\texttt{absolute Schauder frame},\quad \texttt{representing system}.
\]
No prior packet or attempt answered this question. Web searches were also run for close phrases including ``Schauder frame absolute nuclear Frechet'', ``Schauder frames nuclear Frechet absolute'', and ``absolute Schauder frame''. No prior answer was found in the bounded search. Novelty confidence is therefore moderate pending expert review.

\section*{References}

\begin{thebibliography}{9}
\bibitem{CasseseRamos}
G. Cassese and J. P. G. Ramos,
\emph{On the regularity of Fourier interpolation formulas},
arXiv:2503.15733, 2025.
\end{thebibliography}

\end{document}
